\chapter{Expressions}
\label{ch:expressions}

Expressions produce values. Every expression in \haira{} has a type determined at compile time.

\section{Primary Expressions}

\subsection{Literals}

\begin{lstlisting}
42              // int literal
3.14            // float literal
"hello"         // string literal
true            // bool literal
nil             // nil literal
\end{lstlisting}

\subsection{Identifiers}

\begin{lstlisting}
x               // Variable reference
user.name       // Field access
array[0]        // Index access
\end{lstlisting}

\subsection{Parenthesized Expressions}

\begin{lstlisting}
(1 + 2) * 3     // Grouping for precedence
\end{lstlisting}

\section{Arithmetic Expressions}

\begin{lstlisting}
a + b           // Addition
a - b           // Subtraction
a * b           // Multiplication
a / b           // Division
a % b           // Modulo (remainder)
-a              // Unary negation
\end{lstlisting}

\subsection{Integer Division}

Integer division truncates toward zero:

\begin{lstlisting}
7 / 3           // 2
-7 / 3          // -2
7 / -3          // -2
\end{lstlisting}

\subsection{Modulo}

The modulo operator returns the remainder:

\begin{lstlisting}
7 % 3           // 1
-7 % 3          // -1
7 % -3          // 1
\end{lstlisting}

\section{Comparison Expressions}

\begin{lstlisting}
a == b          // Equal
a != b          // Not equal
a < b           // Less than
a > b           // Greater than
a <= b          // Less or equal
a >= b          // Greater or equal
\end{lstlisting}

All comparison expressions produce \type{bool} values.

\subsection{Equality}

Equality (\code{==}) compares by value for primitives and by structural equality for composite types:

\begin{lstlisting}
1 == 1                      // true
"hello" == "hello"          // true
[1, 2] == [1, 2]           // true
User{id: 1} == User{id: 1}  // true (structural)
\end{lstlisting}

\section{Logical Expressions}

\begin{lstlisting}
a and b         // Logical AND (also &&)
a or b          // Logical OR (also ||)
not a           // Logical NOT (also !)
\end{lstlisting}

\subsection{Short-Circuit Evaluation}

Logical operators use short-circuit evaluation:

\begin{lstlisting}
// 'b' is not evaluated if 'a' is false
a and b

// 'b' is not evaluated if 'a' is true
a or b
\end{lstlisting}

\section{String Expressions}

\subsection{Concatenation}

\begin{lstlisting}
"Hello, " + "World!"        // "Hello, World!"
\end{lstlisting}

\subsection{String Interpolation}

\begin{lstlisting}
name = "Alice"
age = 30
"Hello, ${name}! You are ${age} years old."
\end{lstlisting}

Expressions inside \code{\$\{...\}} are evaluated and converted to strings.

\section{Array and Map Expressions}

\subsection{Array Literals}

\begin{lstlisting}
[]                          // Empty array (type from context)
[1, 2, 3]                   // []int
["a", "b", "c"]             // []string
\end{lstlisting}

\subsection{Map Literals}

\begin{lstlisting}
[:]                         // Empty map (type from context)
["a": 1, "b": 2]            // [string:int]
[1: "one", 2: "two"]        // [int:string]
\end{lstlisting}

\subsection{Index Access}

\begin{lstlisting}
array[0]                    // First element
array[array.len - 1]        // Last element
map["key"]                  // Map lookup (returns T?)
\end{lstlisting}

\begin{note}
Array indexing with an out-of-bounds index causes a runtime panic. Map indexing returns an option type.
\end{note}

\section{Field Access}

\begin{lstlisting}
user.name                   // Struct field
point.x                     // Struct field
tuple.0                     // Tuple element
\end{lstlisting}

\section{Function Calls}

\begin{lstlisting}
func()                      // No arguments
func(a, b, c)               // With arguments
obj.method()                // Method call
obj.method(a, b)            // Method with arguments
\end{lstlisting}

\subsection{Chained Calls}

\begin{lstlisting}
text.trim().to_upper().split(" ")
\end{lstlisting}

\section{Pipe Expressions}

The pipe operator passes the left operand as the first argument to the right function:

\begin{lstlisting}
// These are equivalent:
result = f(g(h(x)))
result = x | h | g | f
\end{lstlisting}

\subsection{Pipe with Arguments}

When the right side has arguments, the left value becomes the first argument:

\begin{lstlisting}
// These are equivalent:
result = string.split(text, ",")
result = text | string.split(",")
\end{lstlisting}

\subsection{Complex Pipelines}

\begin{lstlisting}
import "io"
import "string"
import "array"

fn main() {
    data = "  apple, banana, cherry  "

    result = data
        | string.trim
        | string.split(", ")
        | array.map(string.to_upper)
        | array.filter(fn(s) { string.len(s) > 5 })
        | string.join(" | ")

    io.println(result)  // "BANANA | CHERRY"
}
\end{lstlisting}

\section{Range Expressions}

\begin{lstlisting}
0..10           // 0 to 9 (exclusive end)
0..=10          // 0 to 10 (inclusive end)
\end{lstlisting}

Ranges are used primarily in \keyword{for} loops:

\begin{lstlisting}
for i in 0..5 {
    io.println(i)  // 0, 1, 2, 3, 4
}

for i in 0..=5 {
    io.println(i)  // 0, 1, 2, 3, 4, 5
}
\end{lstlisting}

\section{Conditional Expressions}

\subsection{If Expression}

\keyword{if} can be used as an expression:

\begin{lstlisting}
max = if a > b { a } else { b }

status = if score >= 90 {
    "A"
} else if score >= 80 {
    "B"
} else {
    "C"
}
\end{lstlisting}

\begin{warning}
When used as an expression, all branches must have the same type and \keyword{else} is required.
\end{warning}

\subsection{Match Expression}

\keyword{match} expressions provide pattern matching:

\begin{lstlisting}
result = match value {
    0 => "zero"
    1 => "one"
    n if n < 0 => "negative"
    _ => "other"
}
\end{lstlisting}

See Chapter~\ref{ch:statements} for complete pattern matching semantics.

\section{Block Expressions}

Blocks are expressions that evaluate to their last expression:

\begin{lstlisting}
result = {
    x = compute_x()
    y = compute_y()
    x + y  // This is the value of the block
}
\end{lstlisting}

\section{Struct Instantiation}

\begin{lstlisting}
user = User{
    id: 1,
    name: "Alice",
    email: "alice@example.com"
}

// With shorthand (field name matches variable)
name = "Bob"
email = "bob@example.com"
user = User{id: 2, name, email}
\end{lstlisting}

\section{Closure Expressions}

Anonymous functions (closures):

\begin{lstlisting}
// Full syntax
add = fn(a: int, b: int) -> int {
    return a + b
}

// Short syntax for single expression
double = fn(x: int) -> int { x * 2 }

// Type inference for closures
numbers = [1, 2, 3]
doubled = array.map(numbers, fn(n) { n * 2 })
\end{lstlisting}

\section{Option Expressions}

\subsection{Nil Coalescing}

The \keyword{or} operator provides a default for optional values:

\begin{lstlisting}
name = maybe_name or "Unknown"
value = get_value() or default_value
\end{lstlisting}

\subsection{Optional Chaining}

\begin{lstlisting}
// Returns nil if any part is nil
street = user?.address?.street
\end{lstlisting}

\section{Type Assertions}

\begin{lstlisting}
// Assert that value is a specific type
x = value as int

// Safe assertion (returns option)
x = value as? int
\end{lstlisting}

\section{Operator Precedence}

From highest to lowest:

\begin{enumerate}
    \item Postfix: \code{()} \code{[]} \code{.} \code{?.}
    \item Unary: \code{!} \code{-} \code{not}
    \item Multiplicative: \code{*} \code{/} \code{\%}
    \item Additive: \code{+} \code{-}
    \item Range: \code{..} \code{..=}
    \item Pipe: \code{|}
    \item Comparison: \code{<} \code{>} \code{<=} \code{>=}
    \item Equality: \code{==} \code{!=}
    \item Logical AND: \code{and} \code{\&\&}
    \item Logical OR: \code{or} \code{||}
    \item Assignment: \code{=} \code{+=} \code{-=} \code{*=} \code{/=}
\end{enumerate}

\section{Expression Grammar}

\begin{lstlisting}[language=EBNF]
expression = assignment ;

assignment = or_expr [ ("=" | "+=" | "-=" | "*=" | "/=") assignment ] ;

or_expr = and_expr { ("or" | "||") and_expr } ;

and_expr = equality { ("and" | "&&") equality } ;

equality = comparison { ("==" | "!=") comparison } ;

comparison = pipe { ("<" | ">" | "<=" | ">=") pipe } ;

pipe = range { "|" range } ;

range = additive [ (".." | "..=") additive ] ;

additive = multiplicative { ("+" | "-") multiplicative } ;

multiplicative = unary { ("*" | "/" | "%") unary } ;

unary = ("!" | "-" | "not") unary | postfix ;

postfix = primary { call | index | field } ;

call = "(" [ arguments ] ")" ;

index = "[" expression "]" ;

field = "." identifier ;

primary = identifier
        | literal
        | "(" expression ")"
        | array_literal
        | map_literal
        | struct_literal
        | closure
        | if_expr
        | match_expr
        | block ;

arguments = expression { "," expression } ;
\end{lstlisting}
